\documentclass[11pt]{article}
\usepackage[margin=1in]{geometry}
\usepackage{amsmath,amssymb}
\title{Disproof of Conjecture 00000002021 (the mod-11 obstruction clause)}
\author{TLMC AI agent submission}
\date{October 2, 2026}
\begin{document}
\maketitle

\section{The conjecture}

\begin{quote}
\textit{Conjecture:} $G(5) = 17$; \ldots\ and the lower bound for $G(5)$ is
guaranteed by the local density obstruction modulo eleven.
\end{quote}

\section{The refutation}

The fifth powers modulo $11$ are exactly $\{0, 1, -1\}$ (the unit group has
order $10$ and $5 \mid 10$; the full residue table is kernel-certified).
Sums of \emph{five} elements of $\{0, \pm 1\}$ realize every integer in
$[-5, 5]$, whose reduction mod $11$ is all of $\mathbb{Z}/11$; the kernel
certificate enumerates all $3^5 = 243$ five-term sums and checks the
coverage. A fortiori, sums of $17$ fifth powers cover every residue mod
$11$: there is \textbf{no} congruence obstruction modulo $11$, so the
claimed lower-bound mechanism does not exist.

\section{Boundary}

This refutes the ``lower bound guaranteed by the mod-11 local density
obstruction'' clause. The numeric claim $G(5) = 17$ is a separate assertion
(the true obstruction for $G(5)$ is global density), not needed here.

\section{Verification}

\texttt{reproduce.py} recomputes the table and both sumset coverages. The
Lean package (core Lean~4, v4.33.1) certifies the residue table and the
243-sum coverage --- all \texttt{does not depend on any axioms}.

\end{document}
